% ============================================================
%  Artykuł konferencyjny IEEE — szablon dwukolumnowy
%  Tytuł: Zdecentralizowane mechanizmy zapewniania integralności
%         danych modeli 3D z hybrydową walidacją offline
% ============================================================
\documentclass[10pt,conference,a4paper]{IEEEtran}
\IEEEoverridecommandlockouts

\usepackage[utf8]{inputenc}
\usepackage[T1]{fontenc}
\usepackage[polish]{babel}
\usepackage{amsmath}
\usepackage{booktabs}
\usepackage{url}
\usepackage{cite}
\usepackage{graphicx}

% ── Metadane ────────────────────────────────────────────────
\title{Zdecentralizowane mechanizmy zapewniania integralności
       danych modeli 3D z hybrydową walidacją offline}

\author{
  \IEEEauthorblockN{Jakub Kalankiewicz}
  \IEEEauthorblockA{Wydział Elektroniki i Technik Informacyjnych\\
    Politechnika Warszawska\\
    Warszawa, Polska\\
    \texttt{[DO UZUPEŁNIENIA: adres e-mail PW]}}
}

% ============================================================
\begin{document}
\maketitle
% ============================================================

% ── Streszczenie ─────────────────────────────────────────────
\begin{abstract}
Artykuł opisuje architekturę i wyniki badań platformy aukcyjnej
BidSphere, której celem jest bezpieczna dystrybucja modeli~3D
w~formacie GLB. Podstawową luką klasycznego podejścia
klient-serwer jest to, że protokół SSL/TLS chroni wyłącznie
transmisję danych, nie gwarantując integralności po stronie
serwera. Skompromitowany serwer może dostarczyć zmodyfikowany
plik przez poprawne połączenie HTTPS, a odbiorca nie będzie
w~stanie tego wykryć. W~odpowiedzi zaprojektowano dwuwarstwowy
system ochrony łączący uwierzytelniony dostęp przez Cloudinary
z kryptograficzną weryfikacją opartą na blockchainie Ethereum.
Pomiary podstawowe są wykonywane lokalnie, a~wybrane przypadki
zostaną następnie potwierdzone w~sieci testowej Sepolia. Kluczowym wkładem jest implementacja
off-chainowego drzewa Merkle, w~którym na blockchain trafia
wyłącznie korzeń (Merkle Root), co umożliwia całkowicie
offline weryfikację pliku przy użyciu samodzielnego pliku
\texttt{verify.html} bez żadnych wywołań sieciowych.
Przygotowana instrumentacja umożliwia powtarzalne porównanie
czasu zapytania RPC z~lokalnym obliczeniem skrótu i~dowodu
Merkle oraz pomiar kosztu gazu dla różnych rozmiarów partii.
Wartości liczbowe przedstawione w~bieżącej wersji artykułu są
wstępne i~zostaną zastąpione wynikami serii obejmujących trzy
pomiary rozgrzewkowe i~30~powtórzeń.
\end{abstract}

% ── Słowa kluczowe ───────────────────────────────────────────
\begin{IEEEkeywords}
integralność danych, drzewo Merkle, blockchain, weryfikacja
offline, modele 3D, CVE-2012-2459, smart kontrakt, Solidity
\end{IEEEkeywords}

% ── 1. Wstęp ─────────────────────────────────────────────────
\section{Wstęp}

Modele~3D w~formacie GLB zyskują coraz większe znaczenie
w~branżach, gdzie poprawność geometrii jest krytyczna. W~prze-
myśle lotniczym i~wojskowym pliki modeli służą bezpośrednio
do~druku~3D części zamiennych; nawet minimalna zmiana grubości
ścianki może spowodować awarię~\cite{iso_am}. W~medycynie
i~protetyce geometria implantu decyduje o~bezpieczeństwie
pacjenta~\cite{fda_additive}. Na rynku cyfrowych assetów
premium autentyczność pliku bezpośrednio przekłada się na jego
wartość handlową.

Klasyczne podejście opiera się na założeniu, że szyfrowanie
połączenia gwarantuje bezpieczeństwo pliku. Jest to jednak
nieprawdziwe: protokół TLS~1.3 chroni dane \textit{w~locie}
(ang.~\textit{data in transit}), lecz nie zapewnia żadnej
gwarancji integralności po~ich odebraniu ani po~stronie
serwera~\cite{rfc8446}. Atakujący z~dostępem do~infrastruktury
może podmienić plik GLB, a~serwer nadal będzie go dostarczał
przez poprawne połączenie HTTPS. Problem ten określa się jako
brak ochrony danych \textit{w~spoczynku}
(ang.~\textit{data at rest}).

Cel niniejszej pracy jest trójdzielny: (1)~wykrywać podmianę
pliku niezależnie od~stanu serwera, (2)~umożliwiać weryfikację
bez~aktywnego połączenia sieciowego, (3)~minimalizować koszty
operacyjne przez zastosowanie kryptograficznych struktur
off-chain. Praca opisuje platformę BidSphere jako demonstrator
tych właściwości.

% ── 2. Architektura ──────────────────────────────────────────
\section{Architektura Platformy BidSphere\\
         i Wielowarstwowy System Ochrony}

BidSphere zbudowano w~oparciu o~Next.js~16.1.6 z~App~Router, bazę
MongoDB zarządzaną przez ORM~Prisma oraz front-end oparty
na~TailwindCSS. Komunikacja z~blockchainem odbywa się przez
bibliotekę ethers.js~v6; smart kontrakty napisano w~Solidity
0.8.19 i~skompilowano narzędziem Hardhat.

\subsection{Pierwsza warstwa: dostęp i transport}

Uwierzytelnianie użytkowników realizuje NextAuth.js (OAuth~2.0).
Pliki GLB są przechowywane w~Cloudinary jako zasoby chronione
(\textit{authenticated delivery}), do~których dostęp wymaga
podpisanego URL ważnego przez~1~godzinę. Każde żądanie modelu
przechodzi przez serwerowe proxy (\texttt{/api/model/{[}itemId{]}}),
które weryfikuje sesję, generuje podpisany URL i~strumieniuje
odpowiedź z~nagłówkiem \texttt{Cache-Control: no-store},
uniemożliwiając buforowanie po~stronie klienta. Obecna implementacja
rejestruje wizytę na~stronie modelu w~kolekcji \texttt{ModelAccess};
sam endpoint strumieniujący plik nie zapisuje osobnego zdarzenia,
co ogranicza kompletność ścieżki audytu.

\subsection{Druga warstwa: kryptograficzna integralność}

Przy wgrywaniu modelu system oblicza skrót SHA-256 pliku GLB
i~rejestruje go w~smart kontrakcie \texttt{ModelRegistry}
na~sieci Ethereum. Komponent \texttt{VerificationBadge} pobiera
model przez proxy, oblicza skrót po~stronie klienta przy użyciu
\texttt{crypto.subtle.digest()} i~porównuje z~wartością
on-chain. Niezgodność jest natychmiast widoczna jako ostrzeżenie
o~manipulacji. Moduł \texttt{simulateTamper} w~panelu
administratora pozwala na kontrolowane symulowanie ataku przez
podmianę ścieżki do~pliku w~bazie danych przy zachowaniu
oryginalnego skrótu na blockchainie, co umożliwia zebranie
danych eksperymentalnych.

% ── 3. Drzewo Merkle ─────────────────────────────────────────
\section{Zdecentralizowana Optymalizacja\\
         i Walidacja Offline (Drzewo Merkle)}

\subsection{Budowa drzewa i determinizm}

Drzewo Merkle jest binarnym drzewem skrótów, gdzie każdy liść
odpowiada skrótowi SHA-256 jednego modelu. Aby zapewnić
jednoznaczność kolejności wewnątrz każdej pary rodzeństwa,
przed konkatenacją zastosowano sortowanie leksykograficzne:
\begin{equation}
  H_p = \mathrm{SHA\text{-}256}\!\left(
        \min(H_l,H_r)\;\|\;\max(H_l,H_r)\right).
\end{equation}
Drzewo jest dopełniane do~najbliższej potęgi liczby~2 przez
duplikację ostatniego liścia, zgodnie z~podejściem stosowanym
w~protokole Bitcoin~\cite{nakamoto2008}. Implementacja
w~TypeScript (plik \texttt{lib/merkle.ts}) zawiera funkcje
\texttt{buildMerkleTree()}, \texttt{generateProof()}
oraz \texttt{verifyProof()}. Sortowanie par nie uniezależnia
korzenia od~kolejności całej listy liści, dlatego aplikacja
utrzymuje stabilną kolejność modeli w~partii.

\subsection{Smart kontrakt ModelRegistry}

Kontrakt przechowuje dwa typy danych: indywidualne hashe modeli
rejestrowane w~momencie wgrywania pliku oraz korzenie Merkle
reprezentujące całe partie modeli. Kluczowe jest mapowanie O(1)
identyfikujące partię dla~danego modelu:
\begin{verbatim}
mapping(string => uint256)
    private modelToBatchId;
\end{verbatim}
Funkcja \texttt{registerMerkleRoot()} rejestruje korzeń
i~aktualizuje mapowanie dla~wszystkich modeli w~partii
w~jednej transakcji, emitując zdarzenie
\texttt{MerkleRootRegistered}. Kontrakt zapisuje jednak również
tablicę identyfikatorów i~aktualizuje mapowanie dla każdego
modelu, dlatego koszt tej operacji rośnie z~liczebnością partii
i~musi zostać wyznaczony eksperymentalnie.

\subsection{Trójstopniowa weryfikacja offline}

Plik \texttt{verify.html} jest samodzielnym weryfikatorem
niewymagającym żadnych zależności zewnętrznych. Weryfikacja
przebiega w~siedmiu krokach. Krok~1: walidacja struktury
pakietu i~sprawdzenie warunku
\texttt{leafIndex < totalLeaves}. Bez~tej walidacji
sfałszowany wpis z~indeksem w~strefie dopełnienia drzewa
generowałby poprawnie wyglądający dowód Merkle. Jest to ochrona
przed niejednoznacznością analogiczną do problemów z~duplikacją
liści, lecz nie dowód pełnej zgodności modelu ataku z
CVE-2012-2459~\cite{cve2012}. Krok~2:
obliczenie skrótu SHA-256 pliku GLB przez
\texttt{crypto.subtle}. Krok~3: porównanie z~\texttt{fileHash}
z~pakietu dowodowego. Krok~4: iteracyjna rekonstrukcja korzenia
przez haszowanie liścia z~kolejnymi węzłami rodzeństwa.
Krok~5: porównanie ze~skrótem z~pakietu. Krok~6: porównanie
identyfikatora sieci i~adresu kontraktu z~lokalną konfiguracją.
Krok~7: sprawdzenie rekonstruowanego korzenia w~wbudowanej tablicy zaufanych
korzeni. Tablica ta jest kotwicą zaufania (ang.~\textit{trust
anchor}), bez~której atakujący mógłby wygenerować spójny
wewnętrznie, lecz całkowicie sfabrykowany pakiet dowodowy.
Adres kontraktu i~korzenie muszą zostać uzupełnione na podstawie
niezależnie uwierzytelnionego źródła; pusta konfiguracja celowo
uniemożliwia akceptację dowodu. Wyznaczony korzeń jest zawsze
wyświetlany w~interfejsie, umożliwiając późniejsze porównanie
z~Etherscan lub~aneksem pracy dyplomowej. Pola takie jak nazwa,
identyfikator modelu, czas rejestracji, indeks i~liczba liści nie
są częścią haszowanego liścia i~interfejs oznacza je jako
nieuwierzytelnione dane pakietu; indeks i~liczba liści służą
wyłącznie kontroli strukturalnej.

% ── 4. Wyniki ────────────────────────────────────────────────
\section{Analiza Eksperymentalna i Wyniki}

Eksperymenty przeprowadzono przy użyciu modułu
\texttt{simulateTamper}/\texttt{restoreTamper} zaimplementowanego
w~panelu administratora. Tabela~\ref{tab:scenarios} zestawia
cztery scenariusze zabezpieczenia. Czasy weryfikacji mierzono
metodą \texttt{performance.now()} po~stronie klienta w~module
Benchmark; koszty gazu zmierzono narzędziem
\texttt{hardhat-gas-reporter}.

\begin{table}[htbp]
\caption{Wstępne wartości orientacyjne; wyniki końcowe wymagają pełnej serii pomiarowej}
\label{tab:scenarios}
\centering
\renewcommand{\arraystretch}{1.2}
\begin{tabular}{p{1.25cm}p{0.9cm}p{0.95cm}p{1.1cm}p{1.15cm}}
\toprule
\textbf{Scenariusz} & \textbf{Sieć} & \textbf{Wykrycie} &
\textbf{Czas [ms]} & \textbf{Gaz/model} \\
\midrule
Brak ochrony   & Nie  & Nigdy & 0          & 0           \\
Tylko SSL/TLS  & Tak  & Nie   & —          & 0           \\
Blockchain RPC & Tak  & Tak   & $\sim$30$^*$   & $\sim$94\,000$^*$ \\
Merkle offline & \textbf{Nie} & \textbf{Tak} &
                 \textbf{$\sim$5$^*$} & \textbf{zależny od $N$} \\
\bottomrule
\end{tabular}
{\footnotesize $^*$Pojedyncze pomiary wstępne, wyłączone
z~wnioskowania końcowego.}
\end{table}

\subsection{Czas weryfikacji}

Pojedynczy pomiar wstępny dla pliku 5,92\,MB wskazał około
4\,ms dla SHA-256 i~czas poniżej~1\,ms dla przejścia krótkiej
ścieżki Merkle. Lokalny odczyt z~węzła Hardhat trwał około
30\,ms. Wartości te służą wyłącznie do kontroli instrumentacji:
nie pozwalają wyznaczyć rozrzutu ani ekstrapolować opóźnienia
publicznego RPC. Badanie końcowe obejmie dwa profile sieciowe,
trzy rozmiary GLB, trzy pomiary rozgrzewkowe oraz 30 zapisanych
powtórzeń każdej konfiguracji. Raportowane będą mediana,
kwartyle i~p95 oraz pełne dane surowe.

\subsection{Analiza kosztów gazu}

Wstępny raport Hardhat wskazał około 94\,000 jednostek gazu
dla pojedynczego wywołania \texttt{registerModel}; nie jest to
jeszcze estymator wyniku końcowego. Funkcja
\texttt{registerMerkleRoot} zapisuje korzeń i~tablicę
identyfikatorów oraz aktualizuje mapowanie
\texttt{modelToBatchId} dla każdego modelu. Nie można zatem
dzielić kosztu zmierzonego dla jednej partii przez dowolne~$N$
ani traktować kosztu rejestracji korzenia jako stałego.
Dedykowany harness
mierzy rzeczywiste \texttt{gasUsed} dla partii 1, 2, 5, 10,
25, 50 i~100 modeli, porównując sumę~$N$ indywidualnych
transakcji z~jedną transakcją Merkle na świeżo wdrożonych
kontraktach. Wyniki: \textbf{[DO UZUPEŁNIENIA PO SERII 30
POWTÓRZEŃ]}.

% ── 5. Wnioski ───────────────────────────────────────────────
\section{Wnioski i Perspektywy Dalszych Badań}

W~pracy zaprojektowano i~zaimplementowano hybrydowy system
weryfikacji integralności modeli~3D, który adresuje lukę
bezpieczeństwa, której sam protokół SSL/TLS nie jest
w~stanie wypełnić. Płyną z~tego trzy wnioski: (1)~samo
szyfrowanie transportu nie wystarczy dla~danych o~krytycznym
znaczeniu; (2)~kombinacja drzewa Merkle z~blockchainem
umożliwia weryfikację całkowicie offline względem wcześniej
pozyskanego, zaufanego korzenia; (3)~walidacja zakresu indeksu
liścia odrzuca dowody wskazujące zduplikowaną strefę
dopełnienia. Skala oszczędności gazu pozostaje pytaniem
eksperymentalnym, a~dystrybucja korzeni i~klucz właściciela
kontraktu pozostają punktami zaufania.

Pod~względem aplikacyjnym system może znaleźć zastosowanie
w~bazach wojskowych pracujących w~izolowanych sieciach
(weryfikacja geometrii przed~drukiem~3D części zamiennych),
laboratoriach medycznych bez~dostępu do~internetu (walidacja
implantów) oraz~w~systemach zarządzania cyfrowymi assetami
przemysłowymi, gdzie podmiana projektu może nieść
odpowiedzialność prawną.

Głównym kierunkiem dalszych badań jest integracja
z~protokołem \textit{Zero-Knowledge Proofs}~(ZKP), która
pozwoliłaby weryfikować przynależność modelu do~zatwierdzonego
zbioru bez~ujawniania samej geometrii~\cite{boneh2023}.
Innym interesującym rozwinięciem byłoby automatyczne
odświeżanie tablicy zaufanych korzeni w~pliku
\texttt{verify.html} przy~użyciu podpisów cyfrowych
dostarczanych poza pasmem (ang.~\textit{out-of-band}).

% ── Literatura ───────────────────────────────────────────────
\begin{thebibliography}{99}

\bibitem{nakamoto2008}
S.~Nakamoto, ``Bitcoin: A~Peer-to-Peer Electronic Cash
System,'' \textit{bitcoin.org}, 2008. [Online]. Dostępny:
\url{https://bitcoin.org/bitcoin.pdf}

\bibitem{rfc8446}
E.~Rescorla, ``The Transport Layer Security (TLS) Protocol
Version~1.3,'' RFC~8446, Internet Engineering Task Force,
sierpień 2018.

\bibitem{cve2012}
``CVE-2012-2459 Detail,'' National Vulnerability Database,
NIST, 2012.
[Online]. Dostępny:
\url{https://nvd.nist.gov/vuln/detail/CVE-2012-2459}

\bibitem{wood2014}
G.~Wood, ``Ethereum: A~Secure Decentralised Generalised
Transaction Ledger,'' \textit{Ethereum Project Yellow Paper},
2014. [Online]. Dostępny:
\url{https://ethereum.github.io/yellowpaper/paper.pdf}

\bibitem{iso_am}
ISO/ASTM~52900:2021, ``Additive manufacturing --- General
principles --- Fundamentals and vocabulary,'' International
Organization for Standardization, 2021.

\bibitem{fda_additive}
U.S.~Food and Drug Administration, ``Technical Considerations
for Additive Manufactured Medical Devices,'' FDA Guidance
Document, 2017.

\bibitem{boneh2023}
D.~Boneh i~V.~Shoup, ``A~Graduate Course in Applied
Cryptography,'' wersja~0.6, 2023. [Online]. Dostępny:
\url{https://toc.cryptobook.us/}

\bibitem{szabo1997}
N.~Szabo, ``Formalizing and securing relationships on public
networks,'' \textit{First Monday}, t.~2, nr~9, 1997.

\end{thebibliography}

\end{document}
